%%
%% This is file `main.tex' based on `sample-sigconf.tex' (q.v. for spurce of that,
%%
%% IMPORTANT NOTICE:
%% 
%% For the copyright see the original source file `sample-sigconf.tex'
%% in the `Sample' folder.
%%
%% For distribution of the original source see the terms
%% for copying and modification in the file samples.dtx.
%% 
%% This generated file may be distributed as long as the
%% original source files, as listed above, are part of the
%% same distribution. (The sources need not necessarily be
%% in the same archive or directory.)
%%
%% Commands for TeXCount
%TC:macro \cite [option:text,text]
%TC:macro \citep [option:text,text]
%TC:macro \citet [option:text,text]
%TC:envir table 0 1
%TC:envir table* 0 1
%TC:envir tabular [ignore] word
%TC:envir displaymath 0 word
%TC:envir math 0 word
%TC:envir comment 0 0
%%
%%
%% The first command in your LaTeX source must be the \documentclass command.

%% NOTE that a single column version is required for 
%% submission and peer review. This can be done by changing
%% the \doucmentclass[...]{acmart} in this template to 
%%\documentclass[manuscript,review,anonymous]{acmart}
%% This version is used for drafting and final submission
\documentclass[sigconf]{acmart}


%% 
%% To ensure 100% compatibility, please check the white list of
%% approved LaTeX packages to be used with the Master Article Template at
%% https://www.acm.org/publications/taps/whitelist-of-latex-packages 
%% before creating your document. The white list page provides 
%% information on how to submit additional LaTeX packages for 
%% review and adoption.
%% Fonts used in the template cannot be substituted; margin 
%% adjustments are not allowed.

%%
%% \BibTeX command to typeset BibTeX logo in the docs
\AtBeginDocument{%
  \providecommand\BibTeX{{%
    \normalfont B\kern-0.5em{\scshape i\kern-0.25em b}\kern-0.8em\TeX}}}

%% Rights management information.  This information is sent to you
%% when you complete the rights form.  These commands have SAMPLE
%% values in them; it is your responsibility as an author to replace
%% the commands and values with those provided to you when you
%% complete the rights form.
\setcopyright{acmlicensed}
\copyrightyear{2026}
\acmYear{2026}
%%\acmDOI{XXXXXXX.XXXXXXX}

%% These commands are for a PROCEEDINGS abstract or paper.
%%\acmConference[Conference acronym 'XX]{Make sure to enter the correct
 %% conference title from your rights confirmation email}{June 03--05,
 %% 2018}{Woodstock, NY}
%
%  Uncomment \acmBooktitle if th title of the proceedings is different
%  from ``Proceedings of ...''!
%
%\acmBooktitle{Woodstock '18: ACM Symposium on Neural Gaze Detection,
%  June 03--05, 2018, Woodstock, NY} 
%%\acmISBN{978-1-4503-XXXX-X/18/06}


%%
%% Submission ID.
%% Use this when submitting an article to a sponsored event. You'll
%% receive a unique submission ID from the organizers
%% of the event, and this ID should be used as the parameter to this command.
%%\acmSubmissionID{123-A56-BU3}

%%
%% For managing citations, it is recommended to use bibliography
%% files in BibTeX format.
%%
%% You can then either use BibTeX with the ACM-Reference-Format style,
%% or BibLaTeX with the acmnumeric or acmauthoryear sytles, that include
%% support for advanced citation of software artefact from the
%% biblatex-software package, also separately available on CTAN.
%%
%% Look at the sample-*-biblatex.tex files for templates showcasing
%% the biblatex styles.
%%

%%
%% The majority of ACM publications use numbered citations and
%% references.  The command \citestyle{authoryear} switches to the
%% "author year" style.
%%
%% If you are preparing content for an event
%% sponsored by ACM SIGGRAPH, you must use the "author year" style of
%% citations and references.
%% Uncommenting
%% the next command will enable that style.
%%\citestyle{acmauthoryear}

%%
%% end of the preamble, start of the body of the document source.

%% % Location of your graphics files for figures, here a sub-folder to the main project folder
\graphicspath{{./images/}} 


\begin{document}

%%
%% The "title" command has an optional parameter,
%% allowing the author to define a "short title" to be used in page headers.
\title{LLM-Based Anomaly Detection in System Logs}

%%
%% The "author" command and its associated commands are used to define
%% the authors and their affiliations.
%% Of note is the shared affiliation of the first two authors, and the
%% "authornote" and "authornotemark" commands
%% used to denote shared contribution to the research.
\author{Lam Christopher Nguyen}
%%\authornote{Both authors contributed equally to this research.}
\email{la815794@ucf.edu}
\orcid{407-825-6666}
\affiliation{%
  \institution{University of Central Florida}
  \city{Orlando}
  \state{Florida}
  \country{USA}
}


%%
%% By default, the full list of authors will be used in the page
%% headers. Often, this list is too long, and will overlap
%% other information printed in the page headers. This command allows
%% the author to define a more concise list
%% of authors' names for this purpose.
\renewcommand{\shortauthors}{Nguyen}

%%
%% The abstract is a short summary of the work to be presented in the
%% article.
\begin{abstract}

A milestone update will be given on how Large Language Models can be used to identify problems in system and network logs. The practical application that will be used is a web application designed using the MongoDB, Express, React, and Node (MERN) stack. The framework that will be used to scan the network logs is called LangChain combined with Claude. The topic number in this syllabus is: 14. LLM-Based Anomaly Detection in System Logs. So far the application has been built with Backend Logging using the Javascript logging Library Winston which so far is able to log backend API requests and responses. LangChain framework has been implemented to some extent and is able to scan logs. However more fine-tuning needs to be done as it is not consistent and there are still bugs and issues with thee Claude API.

\end{abstract}

%%
%% The code below is generated by the tool at: http://dl.acm.org/ccs.cfm
%% Please copy and paste the code instead of the example below.
%%
\begin{CCSXML}
<ccs2012>
 <concept>
  <concept_id>10010147.10010178.10010179</concept_id>
  <concept_desc>Computing methodologies~Natural language processing</concept_desc>
  <concept_significance>500</concept_significance>
 </concept>
 <concept>
  <concept_id>10002978.10003029.10003032</concept_id>
  <concept_desc>Security and privacy~Intrusion/anomaly detection and malware mitigation</concept_desc>
  <concept_significance>500</concept_significance>
 </concept>
 <concept>
  <concept_id>10010147.10010257.10010293.10010294</concept_id>
  <concept_desc>Computing methodologies~Neural networks</concept_desc>
  <concept_significance>300</concept_significance>
 </concept>
 <concept>
  <concept_id>10011007.10011074.10011111.10011113</concept_id>
  <concept_desc>Software and its engineering~System administration</concept_desc>
  <concept_significance>100</concept_significance>
 </concept>
</ccs2012>
\end{CCSXML}

\ccsdesc[500]{Computing methodologies~Natural language processing}
\ccsdesc[500]{Security and privacy~Intrusion/anomaly detection and malware mitigation}
\ccsdesc[300]{Computing methodologies~Neural networks}
\ccsdesc[100]{Software and its engineering~System administration}

%%
%% Keywords. The author(s) should pick words that accurately describe
%% the work being presented. Separate the keywords with commas.
\keywords{large language models, anomaly detection, log analysis, system monitoring, LangChain, natural language processing, security analysis, network logs}

%%Guidelines for choosing keywords:

%%ACM Digital Library Guide: https://authors.acm.org/proceedings/production-information/preparing-your-article-with-latex
%%General best practices: Choose 5-10 specific terms that researchers would use when searching for your work. Include both broad terms (e.g., "anomaly detection") and specific ones (e.g., "LangChain")
%%Tips for your paper:

%%Use terms from your abstract and title
%%Include the specific technologies/frameworks you're using (LangChain, Claude)
%%Add domain-specific terms (log analysis, system monitoring)
%%Consider what someone would search for to find work like yours
%%Look at keywords in similar papers in your field
\maketitle

\section{Introduction and Problem Statement}


\subsection{Motivation} 

The motivation for this project is to explore how LLMs can be used to identify problems in system and network logs, which can be a valuable tool for system administrators and security professionals. 

\subsection{Significance}

The significance of this project is that it can help to improve the efficiency and effectiveness of log analysis, which can lead to faster identification of issues and improved security. 

\subsection{Objectives}

The objectives of this project are to develop a web application that can scan network logs using LangChain and Claude, and to evaluate the performance of the application in identifying anomalies in the logs.


\section{Related Work}

Recent advances in applying large language models to log analysis have shown promising results. 

Egersdoerfer et al.~\cite{10.1145/3588195.3595943} began using LLMs such as ChatGPT for log analysis since they are able to retain context and formulate insightful responses over conversations. These LLMs are able to conduct few-shot and in-context learning with reasoning ability.
Previously there were several problems with pre-existing techniques. The first issue was that there was a heavy reliance on expert-labeled logs to discern anomalous patterns. This manual labelling slowed down training. The second issues was that they relied on numeric model prediction based on numeric vector input which is not human readable. This lack of human readability rules out targeted error correction by humans. 

Qi et al.~\cite{10466820} introduced LogGPT. This paper explored the use of ChatGPT for log-based anomaly detection. It demonstrated that ChatGPT could be effectively used to identify anomalies in system logs, showing promising results in terms of accuracy and efficiency compared to traditional methods.

Guan et al.~\cite{guan2025logllmlogbasedanomalydetection} introduced LogLLM.Traditional deep-learning methods often struggle to capture semantic information in log messages which are typically organized in natural language. LogLLM employed BERT for extracting semantic vectors from log messages, and uses Llama for classifying log sequences. It also uses a projector to align the vector spaces of BERT and Llama to ensure consistency in understanding log semantics. LogLLM doesn't require log parsers to extraxt templates from logs.

He et al.~\cite{10771398} introduced LLMeLog. This paper introduced a system intended to solve two problems. The first problem is that semantic noise is buried in log events which lead to inevitable gaps between the semantics from the logs and their actual intended meaning. The second problem was that there exists a gap between general semantic embeddings and the specific embedding requirement for anomaly detection tasks. LLMeLog uses LLMs to enrich the content of log events with in-context learning techniques. These enriched logs are then used to fine-tune a pre-trained BERT model. Finally, it trains a transformer-based anomaly detection model with the event representationas produced by this fine-tuned BERT model. F1 scorese exceeded 99 percent and even with using only 10 percent of the labeled data, was still able to achieve an F1 score of 90 percent.


\section{Methodology}

In order to demonstrate how LLMs can be used to identify problems in system and network logs, a web application will be developed using the MERN stack. The application will then utilize some LLM log frameworks that were developed for LLM log analysis to scan a repository of logs and identify anomalies. An attempt will be made to balance performance with cost efficiency. The performance of the application will be evaluated using standard evaluation metrics for anomaly detection.


\subsection{Techniques}

The technique for working with Logging involves 4 steps. 

The first step is to collect and possibly parse the logs. The way the application is designed, all the logs will be generated and downloaded into a single location. The logs will then be parsed and preprocessed to ensure that they are in a format that can be easily analyzed by the LLM.

The second step is to perform Vector Search and Contextual Retrieval. This step provides a baseline of what normal operations shown in the logs will look like. Retrieval augmented generation (RAG) is a technique that can be used to improve the performance of LLMs by providing them with relevant context from the logs. This can be done using vector search techniques to retrieve relevant log entries based on their semantic similarity to the input query. RAG will be used to feed the LLM.The LLM that will be used is Claude.

The third step is to perform Anomaly Detection. This step involves using the LLM to analyze the logs and identify any anomalies or patterns that may indicate a problem. The LLM can be trained to recognize normal patterns in the logs and flag any deviations from those patterns as potential anomalies.

The fourth step is to analyze the root cause of all the errors from the series of logs and explain in human readable language what the problem is and how to fix it. 

\subsection{Evaluation Metrics}

The performance of the application will be evaluated using standard evaluation metrics for anomaly detection, such as precision, recall, and F1 score. These metrics will be calculated based on the number of true positives, false positives, and false negatives identified by the application.

For review, precision is the ratio of true positives to the total number of positive predictions made by the model. Recall is the ratio of true positives to the total number of actual positive cases in the dataset. The F1 score is the harmonic mean of precision and recall, providing a single metric that balances both aspects of performance.

\section{Preliminary Results}

So far the application has been built with Backend Logging using the Javascript logging Library Winston which so far is able to log backend API requests and responses. LangChain framework has been implemented to some extent and is able to scan logs. However more fine-tuning needs to be done as it is not consistent and there are still bugs and issues with thee Claude API. The next steps will be to continue fine-tuning the integration of the LangChain framework and the Claude API to improve function, consistency and reliability. Once the application is able to consistently scan logs and identify anomalies, the next step will be to evaluate its performance using the evaluation metrics mentioned above.


%%
%% The acknowledgments section is defined using the "acks" environment
%% (and NOT an unnumbered section). This ensures the proper
%% identification of the section in the article metadata, and the
%% consistent spelling of the heading.
%%\begin{acks}
%%Acknowledgements go here. Delete enclosing begin/end markers if there are no acknowledgements.
%%\end{acks}

%%
%% The next two lines define the bibliography style to be used, and
%% the bibliography file.
\bibliographystyle{ACM-Reference-Format}
\bibliography{references.bib}

%%
\end{document}
\endinput
%%
%% End of file `main.tex'.
